% =====================================================================
% Section: Method  (outline.md §2)
% Draft v1 (2026-06-10). Faithful to implementation in
%   CATANet/basicsr/archs/catanet_arch.py (DPR / TAB / IASA / CATANet).
% Requires in preamble: amsmath, amssymb, bm (optional).
% Notation is self-contained; symbols are introduced on first use.
% This draft contains NO result assertions (no "results show"); all
% statements are architectural descriptions or design motivations.
% =====================================================================
\section{Method}
\label{sec:method}

\subsection{Overview}
\label{sec:method_overview}
We build on the lightweight super-resolution (SR) backbone of CATANet and replace
its token-routing component with a \emph{Dynamic Prototype Routing} (DPR) module
(Fig.~\ref{fig:arch}).
Given a low-resolution (LR) input $I_{LR}\in\mathbb{R}^{3\times H\times W}$, a
$3\times3$ convolution extracts shallow features
$F_0\in\mathbb{R}^{C\times H\times W}$ ($C{=}40$). The deep feature extractor is a
stack of $L{=}8$ blocks; each block couples a content-routed global aggregation
unit (the Token Aggregation Block, TAB) with a local self-attention unit (LRSA),
followed by a $3\times3$ fusion convolution and a residual connection:
\begin{equation}
F_l = F_{l-1} + \mathrm{Conv}_{3\times3}\!\big(\mathrm{LRSA}_l(\mathrm{TAB}_l(F_{l-1}))\big),
\quad l=1,\dots,L .
\end{equation}
A global residual links the deep features to the network output, which is
reconstructed by a pixel-shuffle upsampler and added to a bilinearly upsampled
base image:
\begin{equation}
I_{SR} = \mathrm{Up}\big(F_0 + F_L\big) + \mathrm{Bilinear}_{\uparrow s}(I_{LR}),
\end{equation}
where $s\in\{2,3,4\}$ is the scale factor. The contribution of this work is
entirely inside TAB: we keep its external interface
$\mathbb{R}^{C\times H\times W}\!\to\!\mathbb{R}^{C\times H\times W}$ unchanged and
substitute the original means-based center iteration and hard-label sorting with
the DPR pipeline (Sec.~\ref{sec:method_dpr}), while reusing the subsequent
intra/inter-group aggregation operator (IASA, Sec.~\ref{sec:method_iasa}). This
keeps the rest of the backbone untouched and isolates the effect of the routing
mechanism.

Inside a TAB, the input feature map is flattened into a token sequence
$X\in\mathbb{R}^{B\times N\times C}$ with $N{=}HW$, layer-normalized, routed by DPR,
aggregated by IASA, and fused by a $1\times1$ convolution and a convolutional
feed-forward network (ConvFFN), each with a residual connection. The data flow is
\begin{equation}
X \;\to\; \mathrm{DPR} \;\to\; \mathrm{IASA} \;\to\; \mathrm{Conv}_{1\times1}
\;\to\; \mathrm{ConvFFN}.
\end{equation}

\subsection{Motivation: limitations of center-based routing}
\label{sec:method_motivation}
The original TAB performs content-aware reordering by maintaining a set of routing
centers via a $k$-means--style iteration with an exponential moving average (EMA)
buffer, assigning each token to its nearest center by an \texttt{argmax} hard
label, sorting tokens by that label, and aggregating the sorted sequence. This
design has four limitations that motivate DPR: (i) the centers are driven by a
cross-batch historical buffer rather than strictly by the current input, so they
adapt slowly to per-image semantics; (ii) the token-to-center relation is a hard
label that discards the assignment confidence; (iii) the ordering only clusters by
label, with no intra-cluster purity constraint; and (iv) prototype generation and
prototype confirmation are entangled, weakening interpretability. DPR addresses
these points by generating prototypes from the current tokens, decoupling
generation from confirmation, and propagating an explicit assignment confidence
through routing and aggregation.

\subsection{Dynamic Prototype Routing}
\label{sec:method_dpr}
DPR maps a token sequence $X\in\mathbb{R}^{B\times N\times C}$ to a content-reordered
sequence together with a set of $M$ dynamic prototypes and routing meta-information.
We describe its four stages (Fig.~\ref{fig:dpr}).

\paragraph{(1) Dynamic prototype generation.}
We first embed tokens with a lightweight head
$E = \mathrm{GELU}(\mathrm{Linear}(\mathrm{LN}(X)))\in\mathbb{R}^{B\times N\times C}$,
and predict a soft assignment of each token to $M$ prototype slots:
\begin{equation}
A = \mathrm{softmax}\big(E W_a\big)\in\mathbb{R}^{B\times N\times M},
\quad W_a\in\mathbb{R}^{C\times M},
\label{eq:assign}
\end{equation}
where the softmax is taken over the $M$ slots. Each prototype is the
assignment-weighted average of the \emph{current} tokens,
\begin{equation}
P^{\mathrm{c}}_{m} = \frac{\sum_{n} A_{n,m}\, X_{n}}{\sum_{n} A_{n,m} + \epsilon},
\qquad
P^{\mathrm{c}} = \mathrm{normalize}\big(\mathrm{LN}(P^{\mathrm{c}})\big),
\label{eq:proto_content}
\end{equation}
with $\ell_2$ normalization along the channel dimension. Unlike the EMA centers of
the baseline, $P^{\mathrm{c}}$ depends only on the current input and carries no
cross-batch state (addressing limitation~(i)).

\paragraph{(2) Prototype confirmation by query refinement.}
To decouple prototype \emph{generation} from \emph{confirmation}
(limitation~(iv)), we refine $P^{\mathrm{c}}$ with a learnable query bank
$Q\in\mathbb{R}^{M\times C}$ via a single cross-attention over the tokens. With
query seed $\tilde{P}=P^{\mathrm{c}}+Q$, queries from the seed, and keys/values
from the embedded/raw tokens,
\begin{equation}
P^{\mathrm{r}} = \mathrm{softmax}\!\Big(\tfrac{(\tilde{P}W_q)(E W_k)^{\!\top}}
{\sqrt{d_r}}\Big)\,(X W_v),
\end{equation}
where $d_r$ is the router dimension. The refinement is fused through a learnable
gate $\gamma=\sigma(g_r)\in(0,1)$,
\begin{equation}
P = \mathrm{normalize}\big(\mathrm{LN}(P^{\mathrm{c}} + \gamma\,P^{\mathrm{r}})\big).
\label{eq:proto_refine}
\end{equation}
Initializing $g_r$ to a small value keeps the refinement a residual perturbation at
the start of training, so the prototype identities remain stable across samples.
This stage can be disabled (ablation A1) by setting
\texttt{use\_prototype\_query\_refine}$=$false, which falls back to
$P=P^{\mathrm{c}}$.

\paragraph{(3) Confidence-aware token-to-prototype matching.}
We match tokens \emph{back} to the confirmed prototypes in a shared
$\ell_2$-normalized router space, $\hat{e}_n=\mathrm{normalize}(E_n W_t)$ and
$\hat{p}_m=\mathrm{normalize}(P_m W_p)$, and compute routing scores with a
\emph{learnable temperature} $\tau$:
\begin{equation}
S = \mathrm{softmax}\big(\tau\,\hat{E}\hat{P}^{\!\top}\big)
\in\mathbb{R}^{B\times N\times M},
\qquad
\tau = \min\!\big(e^{\,\theta},\,\tau_{\max}\big),
\label{eq:scores}
\end{equation}
where $\theta$ is a learnable log-scale clamped by $\tau_{\max}$ for stability.
A larger $\tau$ sharpens the assignment distribution and increases the separation
of routing scores. For each token we read out the hard label and its confidence,
\begin{equation}
x^{\mathrm{s}}_n = \max_m S_{n,m},
\qquad
b_n = \arg\max_m S_{n,m},
\label{eq:belong}
\end{equation}
and \emph{explicitly retain} the confidence $x^{\mathrm{s}}_n$ (addressing
limitation~(ii)).

\paragraph{(4) Confidence-aware ordering.}
Instead of sorting only by the cluster label, we use a composite key that keeps the
inter-cluster order while ordering high-confidence tokens first \emph{within} each
cluster (addressing limitation~(iii)):
\begin{equation}
k_n = b_n + 0.5\,\big(1 - x^{\mathrm{s}}_n\big),
\qquad
\pi = \mathrm{argsort}_n(k_n).
\label{eq:sortkey}
\end{equation}
Since $x^{\mathrm{s}}_n\in(0,1)$, the additive term lies in $(0,0.5)$ and never
crosses an integer label boundary, so clusters stay contiguous while their members
are purity-ordered. The token sequence, the labels, and the confidences are
gathered by $\pi$ into $X^{\pi},\,b^{\pi},\,x^{\mathrm{s},\pi}$, and the inverse
permutation $\pi^{-1}$ is recorded for scatter-back. Setting
\texttt{use\_conf\_sort}$=$false reduces Eq.~\eqref{eq:sortkey} to a pure
label sort $k_n=b_n$ (ablation A2).

DPR outputs the reordered tokens $X^{\pi}$, the dynamic prototypes $P$, the full
score map $S$, and the per-token confidence $x^{\mathrm{s}}$, which are consumed by
the aggregator below.

\subsection{Confidence-aware aggregation (IASA)}
\label{sec:method_iasa}
The aggregation operator is retained from the backbone but is made
confidence-aware. On the reordered sequence $X^{\pi}$ it computes (a) intra-group
local attention over overlapping groups of size $g$, capturing fine local
structure within each contiguous, semantically coherent segment; and (b) a global
branch that attends to the dynamic prototypes $P$ projected to keys/values. Let
$O_{\mathrm{loc}}$ and $O_{\mathrm{glb}}$ denote the two branch outputs. They are
combined through a learnable global gate $\beta=\sigma(g_g)$ modulated by the
per-token confidence:
\begin{equation}
O = O_{\mathrm{loc}} + \beta\cdot x^{\mathrm{s},\pi}\cdot O_{\mathrm{glb}},
\label{eq:iasa_gate}
\end{equation}
so that low-confidence tokens rely less on the global prototype summary. The output
is scattered back to the original token order via $\pi^{-1}$ and projected. Setting
\texttt{use\_iasa\_score\_gate}$=$false removes the confidence factor in
Eq.~\eqref{eq:iasa_gate} ($O = O_{\mathrm{loc}} + \beta\,O_{\mathrm{glb}}$;
ablation A2).

\paragraph{Soft prototype fallback.}
To stabilize tokens whose hard assignment is uncertain, TAB adds a soft, confidence-
weighted prototype context to the aggregated output. With a soft summary
$\bar{X}=S P$ projected by $W_s$ and a learnable gate $\alpha=\sigma(g_s)$,
\begin{equation}
Y = O + \alpha\,\big(1 - x^{\mathrm{s}}\big)\,\big(\bar{X}W_s\big),
\label{eq:soft_fallback}
\end{equation}
where $(1-x^{\mathrm{s}})$ emphasizes low-confidence tokens. The result $Y$ passes
through the $1\times1$ convolution, the residual connection, and the ConvFFN.
Setting \texttt{use\_soft\_fallback}$=$false reduces Eq.~\eqref{eq:soft_fallback}
to $Y=O$ (ablation A2). Equations~\eqref{eq:sortkey}, \eqref{eq:iasa_gate} and
\eqref{eq:soft_fallback} together form the confidence-aware routing path, and the
three switches are independent and default to the full model.

\subsection{Prototype usage balancing}
\label{sec:method_balance}
A learnable temperature can collapse routing onto a few prototypes. To encourage a
balanced use of slots and reduce seed-to-seed variance, we add a usage-entropy
regularizer. Let $u_m=\frac{1}{N}\sum_n S_{n,m}$ be the mean usage of prototype $m$;
its normalized entropy and the auxiliary loss are
\begin{equation}
\mathcal{H}(u) = -\frac{1}{\log M}\sum_{m} u_m\log u_m,
\qquad
\mathcal{L}_{\mathrm{bal}} = \lambda\,\big(1 - \mathcal{H}(u)\big),
\label{eq:balance}
\end{equation}
which is maximized (loss minimized) when usage is uniform. The total objective adds
this term, summed over all blocks, to the pixel reconstruction loss:
\begin{equation}
\mathcal{L} = \mathcal{L}_{\mathrm{pix}} + \sum_{l}\mathcal{L}^{(l)}_{\mathrm{bal}}.
\end{equation}
Setting the per-block weight $\lambda{=}0$ disables balancing (ablation A3).

\subsection{Discussion of design}
\label{sec:method_design}
DPR keeps the backbone's interface and parameter budget in the lightweight regime
while changing only how tokens are organized for aggregation. The four stages map
directly to the contributions: dynamic generation (Eqs.~\ref{eq:assign}--%
\ref{eq:proto_content}) for C1, generation/confirmation decoupling
(Eq.~\ref{eq:proto_refine}) for C2, confidence-aware routing
(Eqs.~\ref{eq:scores}--\ref{eq:sortkey}, \ref{eq:iasa_gate}--\ref{eq:soft_fallback})
for C3, and the learnable temperature with usage balancing (Eqs.~\ref{eq:scores},
\ref{eq:balance}) for C4. Because each design element is exposed as an independent
switch that defaults to the full model and is exactly equivalent to its disabled
form when turned off, the contributions can be isolated cleanly in the ablation
study (Sec.~\ref{sec:ablation}).
